\subsection{A raw two-cut estimate from determinant polynomials}
\label{subsec:binary-raw-surplus}

The two-cut functional has a concrete lower bound on an individual
matching tensor. The estimate below concerns the raw upper functional;
it does not pass either of the two infima in its spectral regularization.
Its operator step is the common-vector lower bound of
\cite[Section~2.4, equation~(41)]{BugarVrana2025}; the determinant
polynomials provide a particular simultaneous Schur-support witness.

\begin{proposition}[A determinant-product witness]
\label{prop:binary-raw-surplus}
Let $g_\alpha$ be the Bug\'ar--Vrana upper functional obtained from
the equal-weight binary geometric mean of the observables for
$12|34$ and $14|23$, where $0<\alpha<1$. For $q\geq2$, put
\[
 W_q=\langle q\rangle_{13}\langle q\rangle_{24},\qquad d=q^2,
 \qquad
 h_d=H\left(\frac{d+1}{2d},\frac1{2d},\ldots,\frac1{2d}\right),
\]
where there are $d-1$ small coordinates and entropy has base two.
Then
\begin{equation}\label{eq:binary-raw-surplus}
 g_\alpha(W_q)\geq2^{(1-\alpha)h_d}.
\end{equation}
In particular,
\[
 \liminf_{\alpha\to0}\log_2g_\alpha(W_2)
 \geq h_4=3-\frac58\log_2 5>1.
\]
\end{proposition}
\begin{proof}
Put \(\mathcal H=(\mathbb C^q)^{\otimes4}\).
Let $D_0$ and $D_1$ be the determinant polynomials of the two
$d$ by $d$ adjacent flattenings. They have degree $d$ and take
values $\pm1$ on $W_q$. In the convention
$D_b(z)=\langle u_b,z^{\otimes d}\rangle$, their symmetric
coefficient vectors have norm one. Indeed, there are $d!$
distinct monomials and $d!$ distinct orderings of each, with
coefficients $\pm1/d!$ in the symmetric tensor.

For $m\geq1$, let $v_m$ represent $(D_0D_1)^m$ and put $N=2dm$.
Orthogonal symmetrization of $u_0^{\otimes m}\otimes u_1^{\otimes m}$
gives
\[
 \|v_m\|\leq1,\qquad
 |\langle v_m,W_q^{\otimes N}\rangle|=1.
\]
For either adjacent cut, multiplication by $D_b^m$ twists both
$\operatorname{GL}_d$ representations by the $m$th determinant
character. Thus each degree-$dm$ Schur partition $\mu$ in the
other factor is shifted to
$\lambda=\mu+(m,\ldots,m)$, with $d$ entries.
Every partition in $v_m$ therefore satisfies $\lambda_i\geq m$
and $|\lambda|=N$. The probability vector defining $h_d$
majorizes every probability vector whose $d$ coordinates are
at least $1/(2d)$. Hence all these Schur entropies are at least
$h_d$, simultaneously for both cuts.

Set $p=(1-\alpha)/\alpha$, and let $A,B$ be the $p$th powers
of the two Schur observables on $\operatorname{Sym}^N(\mathcal H)$.
They are at least $I$. If $P_m$ is the rank-one orthogonal
projection onto $v_m$, the common high-entropy support gives
\[
 A,B\ \geq\ I+(2^{Nph_d}-1)P_m.
\]
Monotonicity and idempotence of the geometric mean imply
$A\#B\geq2^{Nph_d}P_m$. Consequently
\[
 \langle W_q^{\otimes N},(A\#B)W_q^{\otimes N}\rangle
 \geq\frac{2^{Nph_d}}{\|v_m\|^2}
 \geq2^{Nph_d}.
\]
Taking the $\alpha/N$ power and then the defining limit along
$N=2dm$ proves \eqref{eq:binary-raw-surplus}.
\end{proof}

There are two separate limitations. First,
\[
 h_d=\log_2(2d)-\frac{d+1}{2d}\log_2(d+1),
 \qquad h_{4^n}=n+1+o(1).
\]
Since $W_2^{\otimes n}$ is locally equivalent to $W_{2^n}$,
the estimate supplies no fixed positive excess over the tree
threshold per binary matching copy. Power-multiplicativity of
$g_\alpha$ has not been proved. Second, it does not supply the
lower bound uniform in the independent GHZ factor $U_r$.
Thus the strict lower bound for the raw value at $W_2$ alone
does not establish a square price.

\begin{corollary}[A necessary loss under spectral regularization]
\label{cor:binary-regularization-loss}
Assume the price-space theorem of
\cite[Proposition~C.3]{Anthropic2026}, and let
$\widetilde g_\alpha$ be the $\mathbb N$-regularization of the
equal binary-cut functional above. Then
\[
 \limsup_{\alpha\to0}\log_2\widetilde g_\alpha(W_2)\leq\sqrt2,
 \qquad
 \liminf_{\alpha\to0}
 \left(\log_2g_\alpha(W_2)-\log_2\widetilde g_\alpha(W_2)\right)
 \geq h_4-\sqrt2>0.
\]
If the square is not a price, the first upper bound improves to $1$
and the second lower bound improves to $h_4-1$.
\end{corollary}
\begin{proof}
Choose a sequence $\alpha_i\to0$ approaching the upper limit in
the first assertion. Theorem~3.23 of \cite{BugarVrana2025} supplies
dominated normalized LOCC spectral points $f_i\leq g_{\alpha_i}$
with $f_i(W_2)=\widetilde g_{\alpha_i}(W_2)$.
Their scalar orders are $\alpha_i$, by testing reciprocal scalar
tensors. By Lemma~\ref{lem:locc-zero-boundary}, a pointwise subnet
converges to an ordinary normalized restriction spectral point $f$.
The individual-wire formula~\eqref{eq:operator-mean-wire-value}
implies that each cyclic side in its price metric has length at most
$1/2$. Multiplicativity gives the limiting objective $d_{13}+d_{24}$.
The CAT(0) quadrilateral inequality gives
\[
 (d_{13}+d_{24})^2
 \leq2(d_{13}^2+d_{24}^2)
 \leq2(d_{12}^2+d_{23}^2+d_{34}^2+d_{41}^2)
 \leq2.
\]
If the square is not a price, Theorem~\ref{thm:tree-price-dichotomy}
makes this a tree pseudometric, whose four-point condition instead
gives
$d_{13}+d_{24}\leq\max\{d_{12}+d_{34},d_{14}+d_{23}\}\leq1$.
Combine these upper limits with Proposition~\ref{prop:binary-raw-surplus}.
Finally, $5^5<2^{12}$ implies $h_4>3/2>\sqrt2$.
\end{proof}

This loss is forced even without deciding whether the square is a
price. The conclusion does not determine whether the loss occurs
in the tensor-power infimum, in the GHZ-factor infimum, or in both.
